% LTeX: language=es
\section{Teorema de equivalencia de Lax: convergencia por consistencia}

En el ejemplo anterior se aprecia una estructura de prueba muy general.
\eqref{eq:adveccion 1d dependencia continua} nos dice si $U$ es solución y $V$ es ``casi'' solución del esquema, entonces $U$ y $V$ se parecen.
Se puede interpretar como que pequeñas variaciones en la ecuación producen pequeñas variaciones en la solución.
Se dice que el esquema es \emph{estable}.
\begin{definicion}[Esquema estable]

\end{definicion}
Por otro lado \eqref{eq:adveccion 1d consistencia} nos dice que la solución $u$ de \eqref{eq:adveccion} satisface ``razonablemente bien'' el esquema numérico. Se dice que el esquema es \emph{consistente} con la ecuación.
\begin{definicion}[Esquema consistente]

\end{definicion}
Finalmente, se habla de esquemas convergencia
\begin{definicion}[Esquema convergente]

\end{definicion}


\begin{teorema}[Teorema de estabilidad de Lax]
    Un esquema lineal es convergente si y sólo si es estable y consistente.
\end{teorema}

Este es el equivalente en EDPs al teorema de Dahlquist para EDOs.